\section{Marked Type II estimates and rough factors}
\label{sec:type-ii}

The prime construction will count integers whose predecessor has the form
in \Cref{eq:controlled-predecessor}, then remove the composite integers
from that count.  The estimate in this section makes that subtraction
possible: inside a marked bilinear sum, it replaces a prime condition on
one factor by a condition involving only small prime divisors.  We first
state the imported prime estimate and isolate the coefficient property
used by its proof.  We then verify that property for the rough-factor
condition needed when a composite's least prime factor stays below
the square-root range.

Return to the real variable \(x\) from \Cref{sec:reduction}, and put
\begin{equation}\label{eq:sieve-basic-parameters}
 L=\log x,\qquad W=\exp(L^{0.24}),\qquad
 V(y)=\prod_{p\le y}\left(1-\frac1p\right).
\end{equation}
For an integer \(m>1\), write \(P^-(m)\) for its least prime factor, and
put \(P^-(1)=\infty\).  Thus \(P^-(m)>y\) means that \(m\) is
\(y\)-rough.  All Dirichlet characters below, including principal
characters, are extended by zero on nonunits.

Fix an integer \(K\ge1\) and real numbers
\(0.1<a_1<\cdots<a_K<0.2\), all before letting \(x\) grow.  The prime
groups
\begin{equation}\label{eq:marked-prime-groups}
 \mathcal P_i=\{p\text{ prime}:\exp(L^{a_i})\le p\le\exp(2L^{a_i})\}
 \quad (1\le i\le K)
\end{equation}
are disjoint for sufficiently large \(x\).  For a positive integer
\(h\), let
\begin{equation}\label{eq:complete-band-part}
 h_{\mathcal P}=\prod_{p\in\bigcup_i\mathcal P_i}p^{v_p(h)}
\end{equation}
be its complete part supported on these group primes.  A function
\(F:\mathbb Z_{>0}\to\mathbb C\) satisfies \(F(ph)=F(h)\) for every
group prime \(p\) and every positive
\(h\) if and only if it depends only on \(h/h_{\mathcal P}\).  In
particular this invariance removes every power of a group prime, not
only a selected divisor.

Define
\begin{equation}\label{eq:marked-weight}
 \begin{aligned}
 V_i&=\sum_{p\in\mathcal P_i}\frac1p,
 &\omega_i(h)&=\sum_{p\in\mathcal P_i}\one_{p\mid h},\\
 \omega(h)&=\sum_{i=1}^K\omega_i(h),
 &\mathcal W(h)&=q^{\omega(h)-K}
                 \prod_{i=1}^K\frac{\omega_i(h)}{V_i},
                 \qquad 0<q<1.
 \end{aligned}
\end{equation}
The prime number theorem and partial summation give
\(V_i=\log 2+o(1)\).  In particular these denominators are positive
for large \(x\).  The weight is nonnegative and vanishes unless each
group supplies a prime divisor.  It is bounded by a constant depending
only on the fixed \(K,q\), since \(j q^{j-1}\) is bounded for positive
integers \(j\).  In \Cref{sec:prime-construction} we shall take \(q=1/2\).

The following is the one-sided marked Type II estimate of
\emph{The Poisson--Dirichlet law for the prime factors of \(p-1\)}
\cite[Theorem~3.1]{PoissonDirichlet}.
\begin{theorem}[One-sided marked Type II estimate]\label{thm:marked-type-ii}
Fix \(\delta,C,D_*>0\) and \(q\in(0,1)\).  Suppose
\(H_m,H_n\ge x^\delta\) and \(X=H_mH_n\asymp x\).  Let \(F\) satisfy
\(|F(h)|\le1\) and \(F(ph)=F(h)\) for every prime in the groups
defined in \Cref{eq:marked-prime-groups}.  Let \(\alpha_m\) be supported on an
arbitrary interval in \([H_m,2H_m]\), where it has the value
\[
 \alpha_m=m^{iv}\left(\one_{m\text{ prime}}
       -\frac{\one_{P^-(m)>W}}{V(W)\log m}\right),
 \qquad |v|\le L^C.
\]
Let \(\beta_n\) be supported on \(W\)-rough integers in \([H_n,2H_n]\),
with \(|\beta_n|\le L^C\).  If \(K\) is sufficiently large in terms of
\(\delta,C,D_*,q\), then
\begin{equation}\label{eq:marked-type-ii}
 \left|\sum_{m,n}\alpha_m\beta_n F(mn-1)\mathcal W(mn-1)\right|
       \ll X L^{-D_*}.
\end{equation}
The required lower bound on \(K\) is independent of the particular
\(a_i\).  The implicit constant and the threshold for \(x\) may depend
on all the fixed parameters, including the \(a_i\).
\end{theorem}

The invariance is required for every positive \(h\), including when
\(p\mid h\).  The estimate is uniform over the permitted \(F\),
coefficient intervals, and coefficients.  The symbol \(X\) in this
section denotes only the product \(H_mH_n\).

We need the same marked bound for a rough-factor discrepancy.  The
following lemma isolates a sufficient condition on a scalar
coefficient: cancellation against every character of logarithmic
modulus throughout the stated frequency range.  Its proof extracts
the coefficient dependence from the argument of
\cite[Theorem~3.1]{PoissonDirichlet}.

\begin{lemma}[Coefficient criterion for the marked estimate]
\label{lem:marked-coefficient-transfer}
Fix \(\delta,C,D_*>0\) and \(q\in(0,1)\).  Retain the marked data and
the hypotheses on \(F,H_m,H_n\) in \Cref{thm:marked-type-ii}, with fixed
comparison constants in \(X=H_mH_n\asymp x\).  Let
\((\alpha_m)\) and \((\beta_n)\) be scalar sequences indexed by positive
integers, supported on \(W\)-rough integers in their respective dyads,
with \(\alpha\) supported on an arbitrary interval
\(I\subseteq[H_m,2H_m]\), and with
\(|\alpha_m|,|\beta_n|\le L^C\).  Suppose that for every fixed
\(A_0,B,A>0\), every Dirichlet character \(\chi\) modulo
\(k\le L^{A_0}\), and every \(|t|\le2XL^B\),
\begin{equation}\label{eq:type-ii-long-coefficient-interface}
 \left|\frac1{H_m}\sum_m\alpha_m\chi(m)m^{it}\right|\ll_A L^{-A},
\end{equation}
uniformly in the permitted interval, character, and frequency.  The
constant and threshold in this hypothesis may depend on \(A_0,B,A\),
the fixed setup, and all further fixed parameters defining the
coefficients; these parameters are chosen before \(x\) grows.

Then \Cref{eq:marked-type-ii} holds when \(K\) is sufficiently large
in terms of the fixed data, independently of the particular band
exponents \(a_i\).  Its implicit constant and threshold may depend on
all the fixed data, including the \(a_i\).
\end{lemma}

\begin{proof}
We extract the assertion from the proof of
\cite[Theorem~3.1]{PoissonDirichlet}.  In its marked expansion, write
\(mn-1=ah\), where \(a\) is the product of the selected primes.
Invariance gives \(F(mn-1)=F(h)\).  Except when a selected prime has
its square dividing \(mn-1\), one also has
\(\omega(mn-1)-K=\omega(h)\).  Replacing the damping by
\(q^{\omega(h)}\) therefore costs \(O_A(XL^{-A})\) for every fixed
\(A\).  The resulting factor \(F(h)q^{\omega(h)}\) has modulus at
most one and is discarded in the following Cauchy bound.  There are \(O_K(L)\) dyads
\(Y\) for the product of the selected primes.  For the contribution
\(S_Y\) of one dyad, Cauchy's inequality gives
\[
 |S_Y|^2\ll\frac XY\mathcal Q_Y,
\]
where the nonnegative expanded square is
\begin{equation}\label{eq:type-ii-cauchy-square}
 \mathcal Q_Y=\left(\prod_iV_i^{-2}\right)
 \sum_{\substack{a,b,m,n,m',n'\\
        mn-1=ah,\ m'n'-1=bh\text{ for some }h}}
 \rho(a/Y)\rho(b/Y)\alpha_m\overline{\alpha_{m'}}
                         \beta_n\overline{\beta_{n'}}.
\end{equation}
Here \(a,b\) each select one prime from every group, and \(\rho\) is
the fixed real smooth dyadic cutoff supported in \([1,4]\).
Pairs \(a,b\) sharing a selected prime cost \(O_A(XYL^{-A})\).
For the remaining coprime pairs, the common-\(h\) condition is
equivalent to \(t=b-a\), where \(t=bmn-am'n'\).
These reductions use only coefficient bounds and the invariance of
\(F\); see \cite[Section~3.1]{PoissonDirichlet}.

For a fixed \(A_0>0\), let \(\mathfrak M\subset\mathbb R/\mathbb Z\)
be the union of the arcs
\[
 \operatorname{dist}_{\mathbb R/\mathbb Z}(\theta,c/k)
       \le\frac{2L^{A_0}}Y,\qquad
 1\le k\le L^{A_0},\quad c\bmod k,\quad (c,k)=1.
\]
The selected-product range has \(\log Y\gg L^{0.1}\), so these arcs
are disjoint for large \(x\).
Choose the companion's fixed real \(\psi\in C_c^\infty(\mathbb R)\),
equal to one on \([-4,4]\), and define
\[
 \mathcal H_{\mathfrak M}(t;a,b)
 =\psi(t/Y)\int_{\mathfrak M}
       \exp\bigl(2\pi i\theta(t-b+a)\bigr)\,\dd\theta.
\]
The major form is the scalar sum
\begin{equation}\label{eq:type-ii-major-form}
 \begin{split}
 \mathcal Q_Y^{\mathrm{maj}}
 ={}&\left(\prod_iV_i^{-2}\right)
 \sum_{a,b,m,n,m',n'}
 \rho(a/Y)\rho(b/Y)\\
 &\qquad{}\times\alpha_m\overline{\alpha_{m'}}
                         \beta_n\overline{\beta_{n'}}
 \mathcal H_{\mathfrak M}(bmn-am'n';a,b),
 \end{split}
\end{equation}
where \(a,b\) now run independently over products of one prime from
every group.

The proof of the replacement conclusion in
\cite[Proposition~4.1]{PoissonDirichlet} applies to these scalar
coefficients.  Its operator moment is independent of their values.
The passage from that moment to the scalar endpoint pairing uses only
\(W\)-rough support and norm and supremum bounds with logarithmic
exponents depending on \(C\), not on \(K\)
\cite[Sections~3.2--3.3]{PoissonDirichlet}.  This pairing initially
has projections onto the good endpoint states of that argument.  In
removing either projection, only the error from the endpoint outside
its good set is majorized absolutely: each whole scalar endpoint factor
is replaced by its logarithmic supremum bound, discarding both its
coefficient value and its support indicator before the one-edge root
replacement.  The surviving unprojected pairing retains its exact
endpoint vectors.  Expanding those vectors and controlling collisions
uses only size, rough support, and fixed divisor moments
\cite[Section~4.9]{PoissonDirichlet}.  Thus the replacement conclusion,
including restoration of independent \(a,b\), gives
\begin{equation}\label{eq:type-ii-minor-transfer}
 |\mathcal Q_Y-\mathcal Q_Y^{\mathrm{maj}}|
       \ll_A XYL^{-A}
\end{equation}
for every prescribed fixed accuracy, with the choices specified below.

The proof of the major-term estimate
\cite[Proposition~5.1]{PoissonDirichlet} separates characters only at
moduli \(k\le L^{A_0}\).  The group primes and all prime divisors of
the rough variables exceed these moduli, so every variable is a unit
there.  After Mellin separation its retained rectangle, for a fixed
\(B\) chosen after \(A_0\), has \(|s'|\le L^B\) and
\(|t|\le XL^B\), so both \(t\) and \(s'-t\)
lie in the range of \Cref{eq:type-ii-long-coefficient-interface}.
Transform decay and the product mean square on translated intervals
control the discarded tails using only scalar coefficient bounds and
fixed divisor moments.
The further Mellin shift enters only the one-group prime polynomial
and the bounded factor from the other groups; the sparse-frequency
estimate is uniform in this shift
\cite[Lemma~5.3]{PoissonDirichlet}.  On the remaining sparse frequency
set, the only further input about
\(\alpha\) is the supremum in
\Cref{eq:type-ii-long-coefficient-interface}, in the role of
\cite[Lemma~5.2]{PoissonDirichlet}.  The same hypothesis covers the
conjugated endpoint, by conjugating and replacing \(\chi,t\) by
\(\overline\chi,-t\).  The major-term proof therefore gives
\begin{equation}\label{eq:type-ii-major-transfer}
 |\mathcal Q_Y^{\mathrm{maj}}|\ll_D XYL^{-D}
\end{equation}
for every fixed \(D>0\).

Choose the desired square accuracy \(D_1\) sufficiently large in terms
of \(D_*,C\), before choosing \(K\).  For
\Cref{eq:type-ii-minor-transfer}, choose the moment accuracy \(E_0\)
after \(D_1,C\), then the arc exponent \(A_0\), and then \(K\), as in
\cite[Proposition~4.1]{PoissonDirichlet}.  These choices do not depend
on the particular \(a_i\).  The analytic accuracies in
\Cref{eq:type-ii-major-transfer} are chosen after \(A_0\), as permitted
by \Cref{eq:type-ii-long-coefficient-interface} for every fixed
accuracy.  Taking \(A=D=D_1\) gives
\(\mathcal Q_Y\ll XYL^{-D_1}\).  Cauchy's inequality yields
\(|S_Y|\ll XL^{-D_1/2}\).  The \(O_K(L)\) dyads and all fixed
coefficient losses are absorbed by the choice of \(D_1\).  The
logarithmic exponents used in that choice are independent of \(K\);
constants depending on \(K\) affect only the implicit constant and
threshold.  This proves \Cref{eq:marked-type-ii}.
\end{proof}

We now construct a coefficient satisfying this criterion for
\(x^\gamma\)-rough integers.  Its comparison weight comes from the
continuous measure of products of large primes.  For \(w>\gamma>0\),
define
\begin{equation}\label{eq:rough-density-definition}
 \begin{split}
 D_\gamma(w)&=\sum_{j\ge1}D_{\gamma,j}(w),
       \qquad D_{\gamma,1}(w)=\frac1w,\\
 D_{\gamma,j}(w)&=\frac1{j!}
  \int_{\substack{t_i\ge\gamma\ (1\le i<j)\\
                   \sum_{i<j}t_i\le w-\gamma}}
  \frac{1}{w-\sum_{i<j}t_i}\prod_{i<j}\frac{\dd t_i}{t_i}
  \quad(j\ge2).
 \end{split}
\end{equation}
The term with \(j\ge2\) is zero unless \(w\ge j\gamma\), so the sum is
locally finite.  This is the density of
\cite[Lemma~7.4]{BlockPaper}.

Its normalization has an exact continuous meaning.  For an interval
\(I\) whose closure is a compact subset of \((x^\gamma,\infty)\), and
each \(j\ge1\),
\begin{equation}\label{eq:rough-density-pushforward}
 \frac1{j!}\int_{\substack{y_i>x^\gamma\\\prod_i y_i\in I}}
                   \prod_{i=1}^j\frac{\dd y_i}{\log y_i}
 =\frac1L\int_I D_{\gamma,j}(\log y/L)\,\dd y.
\end{equation}
For \(j=1\), the integrands agree because
\(D_{\gamma,1}(\log y/L)/L=1/\log y\).  For \(j\ge2\), put
\(y=\prod_i y_i\), \(t_i=\log y_i/L\) for \(i<j\), and
\(w=\log y/L\).  The change of variables gives
\[
 \prod_{i=1}^j\frac{\dd y_i}{\log y_i}
 =\frac{\dd y}{L}\,
   \frac{\prod_{i<j}(\dd t_i/t_i)}{w-\sum_{i<j}t_i}.
\]
The strict domain from \(y_i>x^\gamma\) differs from the weak simplex
in \Cref{eq:rough-density-definition} only on null boundaries.  The
factor \(1/j!\) divides the ordered prime-product measure by the
number of permutations.  This identity concerns the full domain
\(w>\gamma\); the upper exponent restriction in the next proposition
belongs to the Type II estimate, not to the density.

The functions in \Cref{eq:rough-density-definition} are jointly
continuous on compact subsets of \(w>\gamma>0\), and, for fixed
\(\gamma\), are Lipschitz in \(w\) on each such compact interval.  To
check the assertion also at \(w=j\gamma\) for \(j\ge2\), write
\(s=w-j\gamma\).  For \(s\ge0\), the change of variables
\(t_i=\gamma+s y_i\) gives
\[
 D_{\gamma,j}(w)=\frac{s^{j-1}}{j!}
 \int_{\substack{y_i\ge0\\\sum_{i<j}y_i\le1}}
 \frac{\prod_{i<j}\dd y_i}
 {\prod_{i<j}(\gamma+s y_i)
  \bigl(\gamma+s(1-\sum_{i<j}y_i)\bigr)}.
\]
Every denominator stays bounded away from zero on the indicated compact
sets.  The formula extended by zero for \(s<0\) has the claimed
regularity; the term \(1/w\) is smooth.  In particular \(wD_\gamma(w)\)
has bounded total variation on each fixed compact interval above
\(\gamma\).  This variation bound will allow us to weight the
high-frequency proxy estimate below, and joint continuity will control
the upper Riemann sum in the least-prime-factor subtraction.

The character count proved below gives
\(D_\gamma(\log m/L)/L\) as the principal local density of the
\(x^\gamma\)-rough sequence, while the \(W\)-rough sequence has density
\(V(W)\).  This explains the normalization of its comparison weight.
\begin{proposition}[A rough-factor Type II estimate]\label{prop:rough-type-ii}
Retain the marked data and all the hypotheses on \(F,\beta,H_m,H_n\) of
\Cref{thm:marked-type-ii}.  Fix real numbers
\(0<\gamma<w_-<w_+<1\), and suppose
\[
 [\log H_m/L,\log(2H_m)/L]\subseteq[w_-,w_+].
\]
On an arbitrary interval \(I\subseteq[H_m,2H_m]\), set
\begin{equation}\label{eq:rough-type-ii-coefficient}
 \alpha_m=m^{iv}\left(
    \one_{P^-(m)>x^\gamma}
    -\frac{D_\gamma(\log m/L)}{L V(W)}\one_{P^-(m)>W}
                         \right),\qquad |v|\le L^C,
\end{equation}
and set \(\alpha_m=0\) outside \(I\).  Then
\Cref{eq:marked-type-ii} holds for every fixed \(D_*>0\) when \(K\) is
sufficiently large in terms of the fixed data.  Its required lower bound
is independent of the particular band exponents \(a_i\); the threshold
and implicit constant may depend on all the fixed data.  In particular,
\(\gamma\) and the gap \(w_--\gamma\) are fixed before \(x\) grows.
\end{proposition}

\begin{proof}
We verify \Cref{eq:type-ii-long-coefficient-interface} for this
coefficient.  At low frequencies the two principal counting terms
cancel.  At higher frequencies the prime-product sequence and the
\(W\)-rough comparison sequence are each small.  The support and size
conditions of \Cref{lem:marked-coefficient-transfer} will then follow
directly from the coefficient formula.

Put \(M=H_m\) and write \(R_\gamma(m)=\one_{P^-(m)>x^\gamma}\).
On this dyad a number counted by \(R_\gamma\) has at most
\(J=\lfloor w_+/\gamma\rfloor\) prime factors counted with multiplicity.
Define
\[
 T_\gamma(m)=\sum_{j=1}^{J}\frac1{j!}
       \sum_{\substack{p_1,\ldots,p_j>x^\gamma\text{ prime}\\
                       p_1\cdots p_j=m}}1.
\]
If \(m\) is squarefree and \(x^\gamma\)-rough, the ordered lists give
\(T_\gamma(m)=1\).  If the prime multiplicities are \(e_p\), the
corresponding value is \(1/\prod_p e_p!\), which is at most one.
Thus the difference is supported on a square of a prime exceeding
\(x^\gamma\), and, uniformly for every subinterval of the dyad,
\begin{equation}\label{eq:rough-convolution-square-error}
 \sum_{M\le m\le2M}|R_\gamma(m)-T_\gamma(m)|
 \le\sum_{p>x^\gamma}\left\lfloor\frac{2M}{p^2}\right\rfloor
 \ll Mx^{-\gamma}.
\end{equation}
The same bound holds after any character and real power twist.  It is
smaller than \(M\) times every fixed negative power of \(L\).

We next obtain the counting formula needed at low frequencies.  For
every fixed \(A_1,A_0>0\), every \(k\le L^{A_0}\), every character
\(\chi\bmod k\), and every interval \(I'\subseteq[M,2M]\), we claim
\begin{equation}\label{eq:rough-convolution-low-count}
 \sum_{m\in I'}R_\gamma(m)\chi(m)
 =\one_{\chi\text{ principal}}\frac1L
       \int_{I'}D_\gamma(\log y/L)\dd y+O(M L^{-A_1}).
\end{equation}
The constant may depend on the displayed fixed parameters and on
\(\gamma,w_-,w_+\), but the estimate is uniform in the interval and
character.

To prove it, split every prime variable of a term of \(T_\gamma\) into
half-open dyads \([P_i,2P_i)\).  All their scales satisfy
\(P_i\ge x^\gamma/2\) and \(P_i\le2M\).  There are \(O(L^j)\) possible
boxes for the \(j\)-prime term.  A box that meets the product interval
has
\(M/2^j\le\prod_iP_i\le2M\).  On a prime dyad of scale \(P\), the
Siegel--Walfisz estimate, summed against \(\chi\) over its reduced
classes, gives for every subinterval \(J'\subseteq[P,2P]\)
\begin{equation}\label{eq:rough-slot-siegel-walfisz}
 \sum_{p\in J'}\chi(p)
 =\one_{\chi\text{ principal}}\int_{J'}\frac{\dd y}{\log y}
          +O_{A_2}(P L^{-A_2})
\end{equation}
with arbitrary fixed \(A_2\).  Indeed \(\log P\asymp L\) with fixed
comparison constants, and one chooses the Siegel--Walfisz accuracy to
absorb the at most \(L^{A_0}\) residue classes.  This is the classical
estimate in \cite[Proposition~2.1]{PrimePredecessors}.  Taking
differences of its endpoint estimates gives the absolute error in
\Cref{eq:rough-slot-siegel-walfisz} even for an arbitrarily short
subinterval.  The cutoff \(p>x^\gamma\) can be intersected with that
interval.  For large \(x\), these primes exceed \(k\); hence the
principal character is one on them.  The same orthogonality argument
covers imprimitive characters.

Replace the prime measures in one box successively using
\Cref{eq:rough-slot-siegel-walfisz}.  When all other variables are fixed,
whether at prime values or at real values in a previous replacement, the
product condition leaves an interval in the remaining prime dyad.  Its
error is \(O(P_iL^{-A_2})\).  The total absolute mass of each other
variable, either its prime counting measure or \(\dd y/\log y\), is
\(O(P_h)\).  The error in that box is therefore \(O(M L^{-A_2})\).
There are only the fixed \(j\) replacements and \(O(L^j)\) boxes, so
their errors give \(O(M L^{-A_1})\) on choosing \(A_2\) sufficiently
large.  A nonprincipal character has zero continuous main term.  For a
principal character the main term of the \(j\)-prime summand is exactly
\[
 \frac1{j!}\int_{\substack{y_i>x^\gamma\\\prod_i y_i\in I'}}
                    \prod_{i=1}^j\frac{\dd y_i}{\log y_i}.
\]
By \Cref{eq:rough-density-pushforward}, these integrals sum to the
principal main term in \Cref{eq:rough-convolution-low-count}.
Finally \Cref{eq:rough-convolution-square-error} absorbs repeated
primes and proves the claim.

The matching small-prime count is
\begin{equation}\label{eq:rough-small-prime-character-count}
 \sum_{\substack{m\in I'\\P^-(m)>W}}\chi(m)
 =\one_{\chi\text{ principal}}V(W)|I'|+O(M L^{-A_1})
\end{equation}
for every fixed \(A_1\), with the same interval and character uniformity.
This is \cite[Lemma~2.5]{PoissonDirichlet}.  Notice that all prime divisors
of \(k\le L^{A_0}\) are less than \(W\) for large \(x\), so its principal
character is one on the rough support.

Write \(u=t+v\).  Fix temporarily a constant \(B_2>C+2\).  If
\(|t|\le L^{B_2}\), then \(|u|\le2L^{B_2}\).  Partial summation in
\Cref{eq:rough-convolution-low-count} against \(y^{iu}\) costs at most
\(1+O(|u|)\), since this is its total variation on a dyad.  Partial
summation in \Cref{eq:rough-small-prime-character-count} against
\(D_\gamma(\log y/L)y^{iu}/(LV(W))\) likewise has only a fixed
logarithmic cost: \(D_\gamma\) is Lipschitz on the fixed interval and
Mertens' theorem gives \(LV(W)\asymp L^{0.76}\).  The two resulting
principal main terms are both
\[
 \frac1L\int_I D_\gamma(\log y/L)y^{iu}\dd y,
\]
while both nonprincipal main terms vanish.  Choosing \(A_1\) after
\(B_2\) proves \Cref{eq:type-ii-long-coefficient-interface} with any
fixed accuracy on this low-frequency range.

For the remaining frequencies, use the long prime polynomial estimate
\cite[Lemma~3.1]{PrimePredecessors} and the corresponding
bound for the small-prime proxy.  For fixed
\(0<\tau_{\rm lp}<\theta_{\rm lp}<1\) and fixed \(A_0,A_3>0\), the
prime estimate gives a constant
\(B_0=B_0(\tau_{\rm lp},\theta_{\rm lp},A_0,A_3)\) such that
\begin{equation}\label{eq:rough-slot-high}
 \left|\sum_{p\in J'}\chi(p)p^{-1+iu}\right|\ll L^{-A_3}
\end{equation}
for every interval \(J'\subseteq[P,2P]\), every character modulo
\(k\le L^{A_0}\), and
\[
 x^{\tau_{\rm lp}}/2\le P\le x^{\theta_{\rm lp}},\qquad L^{B_0}\le|u|\le x^2.
\]
Choose \(0<\tau_{\rm lp}<\gamma\) and \(w_+<\theta_{\rm lp}<1\).  Every prime dyad in the
finite convolution lies in this scale range for sufficiently large \(x\).
Partial summation changes the reciprocal weight in
\Cref{eq:rough-slot-high} to give
\(\left|\sum_{p\in J'}\chi(p)p^{iu}\right|\ll P L^{-A_3}\).
In each box of the convolution, fix all but one prime; the remaining
product restriction is again an arbitrary interval.  This bound in that
one slot and trivial bounds in the other slots give \(O(M L^{-A_3})\)
per box.  Taking \(A_3\) to absorb the \(O(L^j)\) boxes and using
\Cref{eq:rough-convolution-square-error} proves any requested
logarithmic saving for
\(M^{-1}\sum_{m\in I}R_\gamma(m)\chi(m)m^{iu}\) throughout this range.

For the proxy, the proof of
\cite[Lemma~5.2]{PoissonDirichlet} gives the following uniform estimate.
Put \(T_* =\exp(L/(\log L)^2)\).  For every interval
\(I'\subseteq[M,2M]\), and with arbitrary fixed \(A_4\), it states
\begin{equation}\label{eq:rough-high-source}
 \begin{split}
 \left|\frac1{M V(W)}
     \sum_{\substack{m\in I'\\P^-(m)>W}}
                \frac{\chi(m)m^{iu}}{\log m}\right|
 \ll L^{-A_4}+{}
 \begin{cases}
 L^{C_0}(|u|^{-1}+x^{-c_0}),&1\le|u|\le T_*,\\
 L^{C_0}\exp(-L/(\log L)^{C_3}),&T_*\le|u|<x^2.
 \end{cases}
 \end{split}
\end{equation}
Here \(C_0,c_0>0\) are fixed after \(A_0,\delta\), and \(C_3>0\) is
absolute; \(C_0\) does not grow with \(A_4\).  The second range uses
the logarithmic-phase estimate of
\cite[Lemma~3.3]{PrimePredecessors}, valid for arbitrary residue
classes and subintervals.  The first follows in the cited proof of
the long-coefficient lemma by comparison with the integral of
\(y^{iu}\) on each progression.

For \(w=\log m/L\), the coefficient of the proxy can be written
\[
 \frac{D_\gamma(w)}{L V(W)}
       =\frac{wD_\gamma(w)}{V(W)\log m}.
\]
The multiplier \(wD_\gamma(w)\) is bounded and has bounded total
variation on our fixed interval.  Since \Cref{eq:rough-high-source}
holds on every subinterval, one more partial summation proves that same
bound for the proxy coefficient.  Choose the requested accuracy first,
then \(A_3,A_4\) and the corresponding \(B_0\), and finally take
\[
 B_2>\max\{C+2,\ B_0+2,\ A+C_0+2\}.
\]
When \(L^{B_2}\le|t|\le2XL^B\), one has
\(|u|\asymp|t|\), \(|u|\ge L^{B_0}\), and
\(|u|\le2XL^B+L^C<x^2\) for large \(x\).  The prime convolution and
\Cref{eq:rough-high-source} therefore give the saving \(L^{-A}\)
for the two parts separately.  The complementary range
\(|t|\le L^{B_2}\), including \(t\) near \(-v\), was already handled
by cancellation of their low-frequency main terms.  This covers every
\(|t|\le2XL^B\); the conjugated assertion then follows as in
\Cref{lem:marked-coefficient-transfer}.

Finally, \(x^\gamma>W\) for large \(x\), and \(D_\gamma\) is bounded
on our fixed interval.  Since \(LV(W)\asymp L^{0.76}\), both parts of
\Cref{eq:rough-type-ii-coefficient} are bounded and supported on
\(W\)-rough integers.  Thus the support, size, and long-coefficient
hypotheses of
\Cref{lem:marked-coefficient-transfer} all hold.  Applying that lemma
proves \Cref{eq:marked-type-ii} and the proposition.
\end{proof}
